% Generated draft: requires mathematical and source review.
\documentclass[11pt]{article}
\usepackage[margin=25mm]{geometry}
\usepackage{fontspec}
\setmainfont{Latin Modern Roman}
\usepackage{amsmath,amssymb,mathtools,unicode-math}
\setmathfont{Latin Modern Math}
\usepackage{enumitem,needspace,xurl,hyperref}
\setlength{\parindent}{0pt}
\setlength{\parskip}{5pt}
\setlength{\emergencystretch}{3em}
\newcommand{\N}{\mathbb{N}}
\newcommand{\Z}{\mathbb{Z}}
\newcommand{\Q}{\mathbb{Q}}
\newcommand{\R}{\mathbb{R}}
\newcommand{\C}{\mathbb{C}}
\newcommand{\F}{\mathbb{F}}
\newcommand{\A}{\mathbb{A}}
\newcommand{\PP}{\mathbb{P}}
\newcommand{\E}{\mathbb{E}}
\newcommand{\eps}{\varepsilon}
\newcommand{\dd}{\,\mathrm{d}}
\begin{document}
\section*{UnsolvedMath ID: 30001451}
\textbf{Source problem code:} catalogue entry
\subsection{Definitions and mathematical statement}
Let \(\mathfrak g\) be a basic classical Lie superalgebra, and let \(\mathcal C\) denote the intended tensor category of \(\mathfrak g\)-supermodules. The supplied record does not specify the ground field, dimensional restrictions on objects, allowed morphisms, or pivotal structure; these must be fixed before the question has a fully precise interpretation.

The notation in the question has the following conventional interpretation, subject to those missing choices. The tensor ideal \(\mathcal I_J\) generated by an object \(J\) consists of direct summands of finite direct sums of objects \(J\otimes X\), with \(X\in\mathcal C\), in the usual symmetric tensor setting for supermodules. A modified trace on this ideal is a family of linear maps
\[
 t_V:\operatorname{End}_{\mathcal C}(V)\longrightarrow k,
 \qquad V\in\mathcal I_J,
\]
where \(k\) is the ground field, satisfying cyclicity and compatibility with the categorical partial trace. Explicitly, for \(f:U\to V\), \(g:V\to U\), and \(h\in\operatorname{End}_{\mathcal C}(U\otimes X)\), these identities are
\[
 t_V(fg)=t_U(gf),\qquad
 t_{U\otimes X}(h)=t_U\bigl(\operatorname{ptr}_X(h)\bigr).
\]
Here \(U,V\in\mathcal I_J\), \(X\in\mathcal C\), and \(\operatorname{ptr}_X\) uses the chosen categorical duality and pivotal structure. An ambidextrous trace on \(J\) is a trace functional satisfying the compatibility between left and right partial traces required for the intended extension to \(\mathcal I_J\). For the resulting family, the modified dimension is conventionally
\[
 d_J(V)=t_V(\operatorname{id}_V).
\]
The record does not supply the precise ambidextrous-trace convention, extension hypotheses, or nonzero normalization. A nonzero scalar rescaling of a chosen trace family does not change whether \(d_J(V)\) vanishes, whereas the zero trace would change the question.

For a simple highest-weight supermodule \(M\) of highest weight \(\lambda\), its usual root-theoretic atypicality \(\operatorname{atyp}(M)\) is the maximum cardinality of a linearly independent, pairwise orthogonal set of isotropic odd positive roots \(\alpha\) satisfying \((\lambda+\rho,\alpha)=0\). Here the positive roots, invariant bilinear form, and \(\rho\), the half-sum of positive even roots minus the half-sum of positive odd roots, belong to the applicable highest-weight setup. The record does not specify that setup or explicitly restrict its modules to a class for which this description applies.

Subject to the intended categorical and highest-weight conventions, the supplied question is: for every basic classical Lie superalgebra \(\mathfrak g\), every simple \(\mathfrak g\)-supermodule \(J\) admitting the intended ambidextrous trace, and every simple \(L\in\mathcal I_J\), does
\[
 \operatorname{atyp}(L)=\operatorname{atyp}(J)
 \quad\Longleftrightarrow\quad
 d_J(L)\ne 0
\]
hold? This is a statement-only formulation of the supplied question; its unresolved ambient hypotheses and trace conventions require source clarification.
\subsection{Short English statement}
The question asks whether, among simple supermodules in the tensor ideal generated by a simple ambidextrous object \(J\), the modified dimension is nonzero exactly for those with the same atypicality as \(J\). The ground field, module category, and trace conventions remain unspecified. The supplied partially solved label is dated curation, not an independently verified mathematical conclusion.
\subsection{Sources}
\begin{itemize}
\item Local UnsolvedMath registry, supplied complete source record: permanent ID \(30001451\), source problem code \texttt{catalogue entry}, title ``Atypicality Detection by Ambidextrous Modified Dimensions.'' The record belongs to the Oberwolfach Reports Open Problems collection and gives the proposed year as 2010; no individual proposer is supplied.
\item Supplied original report citation: \emph{Combinatorial Representation Theory} (2010), Oberwolfach Reports, pp.~799--882, DOI \texttt{10.4171/owr/2010/15}. Original OWR report: \url{https://doi.org/10.4171/owr/2010/15}. The collection describes these reports as documenting mathematical workshops and their open-problem discussions.
\item Supplied statement-verification label: \texttt{corrected\_verified}. The accompanying registry note says that the mathematical content was checked against the official Oberwolfach report and that notation, context, and question form were normalized without changing the source-level claim. This is supplied verification metadata, not a verification performed for this draft.
\item Supplied literature source, labelled ``generalized Kac--Wakimoto for osp'': \url{https://arxiv.org/abs/1112.3384}.
\item Supplied literature source, labelled ``Kac--Wakimoto for the periplectic Lie superalgebra'': \url{https://arxiv.org/abs/1905.04712}.
\item Literature assessment checked at \texttt{2026-08-21}: the supplied status is partially solved. The supplied assessment states that the ordinary Kac--Wakimoto conjecture and generalized versions are proved in important cases, including basic classical, general-linear, or orthosymplectic settings, while the broad generalized ambidextrous-trace formulation in the extracted record is not fully settled. This is supplied dated curation derived from the OWR collection, not independently verified here; its time-sensitive status and cited claims require independent verification.
\end{itemize}
\end{document}
